% !TEX program = pdflatex
\documentclass[12pt]{article}

\usepackage{sbc-template}
\usepackage{graphicx}
\usepackage[hyphens]{url}
\urlstyle{same}
\usepackage[utf8]{inputenc}
\usepackage[brazilian]{babel}
\usepackage{amsmath}
\usepackage{booktabs}
\usepackage{subcaption}

\sloppy

% Ajustes para o LaTeX aproveitar melhor o espaço com figuras
\renewcommand{\topfraction}{0.9}
\renewcommand{\bottomfraction}{0.8}
\renewcommand{\textfraction}{0.1}
\renewcommand{\floatpagefraction}{0.8}
\setlength{\textfloatsep}{8pt plus 2pt minus 2pt}
\setlength{\floatsep}{8pt plus 2pt minus 2pt}
\setlength{\intextsep}{8pt plus 2pt minus 2pt}

\title{Implementação e Comparação de Algoritmos de Ordenação}

\author{Antônio Gemesson Sousa de Oliveira\inst{1}, Gleidson Luan Sena Alves\inst{1}, \\ Ivan Vitor Dias de Oliveira\inst{1}}

\address{Departamento de Computação -- Universidade Federal do Piauí (UFPI)\\
  Teresina -- PI -- Brasil
  \email{\{antonio.gemesson, gleidson.alves, ivan.oliveira\}@ufpi.edu.br}
}

\begin{document}

\maketitle

\begin{abstract}
  This paper presents the implementation and experimental comparison of four
  sorting algorithms: Insertion Sort, Merge Sort, Counting Sort and Radix Sort.
  The algorithms are evaluated over random integer arrays of varying sizes and
  value ranges, relating experimental results to their theoretical complexity
  analyses. Counting Sort was the fastest when the value range is close to the
  input size ($k \approx n$), but it degrades as $k$ grows; Radix Sort was the
  most robust to the value range, and Insertion Sort proved unfeasible for
  $n = 10^5$.
\end{abstract}

\begin{resumo}
  Este artigo apresenta a implementação e a comparação experimental de quatro
  algoritmos de ordenação: Insertion Sort, Merge Sort, Counting Sort e
  Radix Sort. Os algoritmos são avaliados sobre vetores de inteiros aleatórios
  com diferentes tamanhos e intervalos de valores, relacionando os resultados
  experimentais às respectivas análises de complexidade teórica. O Counting
  Sort foi o mais rápido quando o intervalo de valores é próximo do tamanho do
  vetor ($k \approx n$), mas degrada-se à medida que $k$ cresce; o Radix Sort
  foi o mais robusto em relação a $k$, e o Insertion Sort mostrou-se inviável
  para $n = 10^5$.
\end{resumo}

% ============================================================
\section{Introdução}
% ============================================================

A ordenação de dados é um dos problemas fundamentais da ciência da
computação. Ela está presente em tarefas como a organização de registros em
bancos de dados, a busca eficiente em estruturas ordenadas e a preparação de
dados para outros algoritmos. Por isso, a escolha do algoritmo de ordenação
mais adequado pode ter impacto direto no desempenho de uma aplicação.

Os algoritmos de ordenação podem ser divididos em duas grandes classes. Os
algoritmos baseados em comparação, como o Insertion Sort~\cite{gfg-insertion}
e o Merge Sort~\cite{gfg-merge}, decidem a posição dos elementos comparando-os
entre si e, por isso, estão sujeitos ao limite inferior de $\Omega(n \log n)$
comparações no pior caso~\cite{cormen2012}. Já os algoritmos lineares, como o
Counting Sort~\cite{gfg-counting} e o Radix Sort~\cite{gfg-radix}, não
comparam elementos: eles exploram a estrutura numérica das chaves, o que
permite tempo linear sob certas condições, em geral relacionadas ao intervalo
de valores a serem ordenados.

Neste contexto, o objetivo deste trabalho é implementar e comparar
experimentalmente quatro algoritmos de ordenação (Insertion Sort, Merge
Sort, Counting Sort e Radix Sort), relacionando o desempenho observado às
respectivas análises de complexidade teórica. Os algoritmos são avaliados
sobre vetores de inteiros aleatórios, variando tanto o tamanho $n$ da entrada
quanto o intervalo de valores $k$, a fim de investigar em que condições cada
abordagem é mais vantajosa.

% ============================================================
\section{Descrição e Análise Teórica}
\label{sec:teoria}
% ============================================================

\subsection{Insertion Sort}

O Insertion Sort é um algoritmo de ordenação por comparação que constrói a
sequência ordenada de forma incremental, inserindo cada novo elemento na
posição correta em relação aos elementos já ordenados. O algoritmo mantém um
subvetor ordenado à esquerda e, a cada iteração, retira o próximo elemento
não ordenado e o desloca para a posição correta dentro desse subvetor,
movendo os elementos maiores uma posição à frente~\cite{gfg-insertion}.

\textbf{Análise de complexidade:}
\begin{itemize}
  \item \textbf{Melhor caso -- $O(n)$:} Ocorre quando o vetor já está
    ordenado. Em cada iteração, o elemento corrente é maior ou igual ao
    último elemento do subvetor ordenado, e nenhuma troca é realizada. O
    algoritmo executa apenas as $n - 1$ comparações do laço externo.
  \item \textbf{Caso médio -- $O(n^2)$:} Para entradas com distribuição
    aleatória, espera-se que, em média, cada elemento precise ser deslocado
    metade do subvetor ordenado à sua esquerda. O número de comparações e
    deslocamentos cresce quadraticamente com $n$.
  \item \textbf{Pior caso -- $O(n^2)$:} Ocorre quando o vetor está ordenado
    em ordem decrescente. Cada elemento inserido precisa ser comparado e
    deslocado em relação a todos os elementos já ordenados, resultando em
    $O(n^2)$ operações.
\end{itemize}

\subsection{Merge Sort}

O Merge Sort é um algoritmo de ordenação por comparação baseado na estratégia
de divisão e conquista. O vetor é dividido recursivamente ao meio até que cada
subvetor contenha um único elemento (trivialmente ordenado). Em seguida, os
subvetores são mesclados dois a dois em ordem crescente, produzindo subvetores
cada vez maiores até que o vetor completo esteja ordenado. A etapa de
mesclagem percorre os dois subvetores simultaneamente, copiando sempre o
menor elemento para o vetor resultante~\cite{gfg-merge}.

\textbf{Análise de complexidade:}
\begin{itemize}
  \item \textbf{Melhor, médio e pior caso -- $O(n \log n)$:}
    Independentemente da distribuição inicial dos dados, o Merge Sort sempre
    divide o vetor em duas metades e mescla os resultados. A árvore de
    recursão possui altura $O(\log n)$ e, em cada nível, o trabalho total
    de mesclagem é $O(n)$. Portanto, a complexidade é $O(n \log n)$ em
    todos os casos.
  \item \textbf{Espaço auxiliar -- $O(n)$:} A etapa de mesclagem requer
    vetores auxiliares para armazenar temporariamente os elementos,
    resultando em uso de memória proporcional ao tamanho da entrada.
\end{itemize}

\subsection{Counting Sort}

O Counting Sort é um algoritmo de ordenação linear que não realiza
comparações entre elementos. Ele opera contando o número de ocorrências de
cada valor distinto no vetor de entrada e, a partir de uma soma de prefixo
sobre o vetor de contagem, determina a posição exata de cada elemento no
vetor de saída. A iteração do vetor de saída de trás para frente garante a
estabilidade do algoritmo~\cite{gfg-counting}.

\textbf{Análise de complexidade:}
\begin{itemize}
  \item \textbf{Todos os casos -- $O(n + k)$:} onde $n$ é o número de
    elementos e $k$ é o valor do maior elemento, isto é, o limite superior
    do intervalo de valores (ver Seção~\ref{sec:metodologia}).
    O algoritmo realiza uma passada de $O(n)$ para encontrar o máximo, uma
    de $O(n)$ para contar as ocorrências, uma de $O(k)$ para calcular a
    soma de prefixo e uma de $O(n)$ para reconstruir o vetor de saída.
\end{itemize}

\textbf{Parâmetros além do tamanho do vetor:}
Ao contrário dos algoritmos de comparação, a complexidade do Counting Sort
depende não apenas de $n$, mas também de $k$. Quando $k = O(n)$, o algoritmo
é efetivamente linear. Porém, quando $k$ cresce muito além de $n$ (como nos
intervalos $[1, 100n]$ e $[1, 1000n]$ previstos nos experimentos), o custo de
$O(k)$ passa a dominar, degradando o desempenho. Além disso, o Counting Sort
requer $O(k)$ de memória auxiliar, o que pode ser proibitivo para valores de
$k$ muito elevados.

\subsection{Radix Sort}

O Radix Sort é um algoritmo de ordenação linear que ordena os elementos
dígito a dígito, do menos significativo ao mais significativo. Em cada
passada, os elementos são distribuídos com base no dígito corrente usando
uma versão estável do Counting Sort (com soma de prefixo), preservando a
ordem relativa de elementos com o mesmo dígito. A estabilidade é essencial
para que as passadas anteriores não sejam desfeitas. O dígito na posição
\texttt{exp} é extraído por \texttt{(arr[i] / exp) \% 10}~\cite{gfg-radix}.

\textbf{Análise de complexidade:}
\begin{itemize}
  \item \textbf{Todos os casos -- $O(d \cdot n)$:} onde $d$ é o número de
    dígitos do maior elemento e $n$ é o número de elementos. Cada uma das
    $d$ passadas executa em $O(n)$, pois o Counting Sort sobre dígitos de
    0 a 9 usa um vetor de contagem de tamanho fixo (10 posições). Como
    $d = O(\log k)$, a complexidade total é $O(n \log k)$.
  \item \textbf{Vantagem sobre o Counting Sort:} O Radix Sort é menos
    sensível à magnitude dos valores, pois o custo por passada depende
    apenas de $n$ e não de $k$ diretamente.
\end{itemize}

% ============================================================
\section{Metodologia}
\label{sec:metodologia}
% ============================================================

Os algoritmos foram implementados em C++ e operam sobre vetores de inteiros
(\texttt{std::vector<int>}). Os experimentos foram conduzidos nas seguintes
condições:

\begin{itemize}
  \item \textbf{Linguagem de programação e compilador:} C++ e gcc
    version 15.3.1 (Red Hat 15.3.1-1), compilado sem flags de otimização
    (nível padrão do gcc, equivalente a \texttt{-O0}).
  \item \textbf{Sistema operacional:} Linux Fedora.
  \item \textbf{Processador:} Intel Core i5-13420H de 13ª geração.
  \item \textbf{Memória RAM:} 8\,GB DDR5 5200\,MHz.
\end{itemize}

Foram testados os seguintes tamanhos de entrada $n$: $1\,000$, $5\,000$,
$10\,000$, $50\,000$ e $100\,000$. Para cada $n$, os valores dos vetores foram
sorteados com distribuição uniforme discreta, usando o gerador
\texttt{std::mt19937} inicializado com \texttt{std::random\_device}.

Denota-se por $k$ o limite superior do intervalo de valores: os elementos são
sorteados em $[1, k]$, e os quatro intervalos testados correspondem a
$k = n$, $k = 10n$, $k = 100n$ e $k = 1000n$. Esse parâmetro afeta
diretamente o Counting Sort e, indiretamente, o Radix Sort, por meio do
número de dígitos $d$.

Para cada combinação de $n$ e $k$, foram gerados 5 vetores distintos, gravados
em arquivo, de modo que os quatro algoritmos ordenam exatamente a mesma
entrada. Cada algoritmo foi executado 5 vezes por vetor (25 execuções por
combinação), lendo o vetor do arquivo a cada execução, e o tempo reportado
corresponde à média dessas 25 execuções. Como a semente do gerador é obtida
de \texttt{std::random\_device}, os vetores gerados não são reproduzíveis entre
execuções do programa.

O tempo foi medido com \texttt{std::chrono::high\_resolution\_clock}, imediatamente
antes e depois da chamada da função de ordenação. Assim, a leitura do arquivo
e a gravação do resultado não entram na medição, mas a alocação de memória
auxiliar feita pelo próprio algoritmo (como o vetor de contagem do Counting
Sort) é contabilizada. Os algoritmos foram avaliados sequencialmente, sem
paralelismo.

% ============================================================
\section{Resultados e Discussão}
\label{sec:resultados}
% ============================================================

Os resultados são apresentados em quatro etapas: panorama em escala log-log
(Seção~\ref{subsec:loglog}), curvas em escala linear
(Seção~\ref{subsec:linear}), sensibilidade ao intervalo de valores~$k$
(Seção~\ref{subsec:range}) e mapa de calor do \textit{Counting Sort}
(Seção~\ref{subsec:heatmap}).

% -------------------------------------------------------------------
\subsection{Panorama Global: Crescimento Assintótico em Escala Log-Log}
\label{subsec:loglog}
% -------------------------------------------------------------------

As Figuras~\ref{fig:loglog_kn} e~\ref{fig:loglog_media} comparam o tempo
médio dos quatro algoritmos em escala log-log, para $n$ de $10^{3}$ a
$10^{5}$, em dois cenários quanto à relação entre $n$ e $k$. A escala
log-log acomoda a variação de quase seis ordens de grandeza no tempo (de
${\approx}\,2{,}5 \times 10^{-5}$\,s a ${\approx}\,19{,}7$\,s).

\subsubsection{Cenário 1: $k = n$}

Quando $k = n$, cada algoritmo opera em seu regime teórico mais
representativo, e uma complexidade $T(n) = \Theta(n^{\alpha})$ aparece como
uma reta de inclinação~$\alpha$ (Figura~\ref{fig:loglog_kn}).

\begin{figure}[!htbp]
  \centering
  \includegraphics[width=0.75\columnwidth]{Images/comparativo_crescimento.png}
  \caption{Tempo médio de execução em escala log-log com $k = n$, para
           $n \in \{10^3, 5 \times 10^3, 10^4, 5 \times 10^4, 10^5\}$.}
  \label{fig:loglog_kn}
\end{figure}

% TODO (opcional): informar as inclinações medidas (regressão linear em
% log-log) no lugar de "próxima de 2,0" e "próxima de 1,0".
Os resultados alinham-se com a teoria: o \textit{Insertion Sort} tem
inclinação próxima de $2{,}0$ ($\Theta(n^{2})$); o \textit{Merge Sort},
ligeiramente superior a $1{,}0$ ($\Theta(n \log n)$); e o
\textit{Counting Sort} e o \textit{Radix Sort}, próxima de $1{,}0$
(linearidade quando $k = \mathcal{O}(n)$). Para $n = 10^{5}$, a ordem é:
\textit{Counting Sort} (${\approx}\,2{,}3 \times 10^{-3}$\,s),
\textit{Radix Sort} (${\approx}\,8{,}7 \times 10^{-3}$\,s),
\textit{Merge Sort} (${\approx}\,4{,}0 \times 10^{-2}$\,s) e
\textit{Insertion Sort} (${\approx}\,19{,}7$\,s), uma diferença de
aproximadamente quatro ordens de grandeza entre o mais rápido e o mais lento.
O \textit{Radix Sort} fica entre o \textit{Counting Sort} e o \textit{Merge
Sort} porque, com $k = n$, o número de dígitos $d = \mathcal{O}(\log n)$
deixa de ser desprezível, afastando-o do \textit{Counting Sort} à medida que
$n$ cresce.

\subsubsection{Cenário 2: média sobre todos os intervalos de $k$}

A Figura~\ref{fig:loglog_media} apresenta, para cada $n$, a média do tempo
sobre os quatro intervalos testados ($[1, n]$ a $[1, 1000n]$), agregando
desde condições favoráveis ($k = n$) até adversas ($k = 1000n$), o que se
aproxima de aplicações reais, em que o intervalo dos dados raramente é
controlado.

\begin{figure}[!htbp]
  \centering
  \includegraphics[width=0.75\columnwidth]{Images/comparativo_media_geral_n.png}
  \caption{Tempo médio geral de execução em escala log-log, calculado
           como a média sobre todos os intervalos de valores testados
           para cada~$n$. A elevação da curva do \textit{Counting Sort}
           evidencia sua sensibilidade ao parâmetro~$k$.}
  \label{fig:loglog_media}
\end{figure}

Os algoritmos baseados em comparação (\textit{Insertion Sort} e
\textit{Merge Sort}) têm curvas inalteradas entre os cenários, pois suas
complexidades dependem apenas de~$n$. O impacto é drástico para o
\textit{Counting Sort}: sua curva sobe aproximadamente duas ordens de
grandeza (de ${\approx}\,2{,}5 \times 10^{-5}$\,s para
${\approx}\,1{,}8 \times 10^{-3}$\,s com $n = 10^{3}$, e de
${\approx}\,2{,}3 \times 10^{-3}$\,s para ${\approx}\,0{,}20$\,s com
$n = 10^{5}$), e ele ultrapassa o \textit{Merge Sort}, tornando-se o segundo
mais lento, abaixo apenas do \textit{Insertion Sort}. A inversão ocorre
porque os intervalos $k = 100n$ e $k = 1000n$ dominam a média: o custo
$\mathcal{O}(k)$ de inicializar e varrer o vetor de contagens cresce
desproporcionalmente ao trabalho útil sobre os $n$ elementos.

O \textit{Radix Sort} sofre uma elevação muito mais modesta: sua
sensibilidade a~$k$ é logarítmica, e aumentar~$k$ de $n$ para $1000n$
acrescenta apenas $\lceil \log_{10}(1000) \rceil = 3$ passadas. Com isso, ele
passa a ser o algoritmo mais rápido no Cenário~2 (era o segundo no
Cenário~1) e mostra-se a alternativa mais robusta entre os lineares quando o
intervalo não é controlado, enquanto a vantagem do \textit{Counting Sort} é
máxima apenas quando $k \approx n$.

% -------------------------------------------------------------------
\subsection{Curvas de Tempo em Escala Linear}
\label{subsec:linear}
% -------------------------------------------------------------------

A escala log-log comprime diferenças absolutas, por isso as curvas foram
também analisadas em escala linear, separando os algoritmos em dois grupos.

\subsubsection{Insertion Sort: a inviabilidade do crescimento quadrático}

A Figura~\ref{fig:linear_insertion} mostra o \textit{Insertion Sort} em
escala linear, nos dois cenários. A curva é claramente parabólica: os pontos para
$n \leq 10^{4}$ mal se distinguem do eixo horizontal, e o tempo cresce
rapidamente a partir de $n \approx 5 \times 10^{4}$, atingindo
${\approx}\,19{,}7$\,s em $n = 10^{5}$. As curvas dos dois painéis (com
$k = n$ e na média sobre os intervalos de~$k$) são visualmente idênticas, o
que ilustra a independência de~$k$ nos algoritmos de comparação. A classe $\mathcal{O}(n^{2})$ torna-se proibitiva
muito antes de entradas de grande escala, justificando-se apenas para $n$
pequeno ou vetores quase ordenados ($\mathcal{O}(n)$).

\begin{figure}[!htbp]
  \centering
  \begin{subfigure}[t]{0.495\columnwidth}
    \includegraphics[width=\linewidth]{Images/escala_real_insertion_sort.png}
    \caption{Tempo com $k = n$.}
    \label{fig:linear_insertion_kn}
  \end{subfigure}
  \hfill
  \begin{subfigure}[t]{0.495\columnwidth}
    \includegraphics[width=\linewidth]{Images/escala_real_insertion_sort_media_ranges.png}
    \caption{Tempo médio sobre todos os intervalos de $k$.}
    \label{fig:linear_insertion_media}
  \end{subfigure}
  \caption{Tempo de execução do \textit{Insertion Sort} em escala linear.
           O crescimento parabólico confirma a complexidade
           $\mathcal{O}(n^{2})$, e a identidade visual entre os dois
           painéis evidencia a independência do algoritmo em relação
           a~$k$.}
  \label{fig:linear_insertion}
\end{figure}

\subsubsection{Merge Sort, Radix Sort e Counting Sort: eficiência em
escala}

A Figura~\ref{fig:linear_eficientes} agrupa os outros três algoritmos. No
cenário $k = n$ (painel esquerdo), todos ficam abaixo de $0{,}05$\,s para
$n = 10^{5}$, cerca de 500 vezes menos que o \textit{Insertion Sort}, e o
\textit{Counting Sort} tem o menor custo absoluto, pois realiza apenas quatro
varreduras lineares e sua constante oculta em $\mathcal{O}(2n)$ é baixa.

\begin{figure}[!htbp]
  \centering
  \begin{subfigure}[t]{0.495\columnwidth}
    \includegraphics[width=\linewidth]{Images/escala_real_nlog_lineares.png}
    \caption{Tempo com $k = n$.}
    \label{fig:linear_eficientes_kn}
  \end{subfigure}
  \hfill
  \begin{subfigure}[t]{0.495\columnwidth}
    \includegraphics[width=\linewidth]{Images/escala_real_nlog_lineares_media_ranges.png}
    \caption{Tempo médio sobre todos os intervalos de $k$.}
    \label{fig:linear_eficientes_media}
  \end{subfigure}
  \caption{Tempo de execução em escala linear para \textit{Merge Sort},
           \textit{Radix Sort} e \textit{Counting Sort}. À direita, o
           \textit{Counting Sort} passa a ser o mais lento dos três.}
  \label{fig:linear_eficientes}
\end{figure}

No painel direito (média sobre todos os intervalos de~$k$), o
\textit{Counting Sort} sobe para ${\approx}\,0{,}20$\,s em $n = 10^{5}$,
superando o \textit{Merge Sort} (${\approx}\,0{,}040$\,s) por um fator
de~5 e tornando-se o mais lento dos três. O \textit{Radix Sort} é o mais
rápido (${\approx}\,1{,}1 \times 10^{-2}$\,s), contra
${\approx}\,8{,}7 \times 10^{-3}$\,s quando $k = n$, onde é o segundo mais
rápido; essa pequena variação confirma sua robustez frente a~$k$.

% -------------------------------------------------------------------
\subsection{Sensibilidade ao Intervalo de Valores ($k$)}
\label{subsec:range}
% -------------------------------------------------------------------

A Figura~\ref{fig:range} mostra o tempo médio de execução em função
de~$k$, de $10^{3}$ a $10^{6}$, para $n = 1\,000$ fixo.

\begin{figure}[!htbp]
  \centering
  \includegraphics[width=0.75\columnwidth]{Images/impacto_range_n_1000.png}
  \caption{Tempo médio de execução em função do intervalo de valores~$k$
           para $n = 1\,000$ fixo (escala log-log). O \textit{Counting
           Sort} supera o \textit{Radix Sort} já em $k = 10^{4}$, o
           \textit{Merge Sort} entre $k = 10^{4}$ e $k = 10^{5}$, e o
           \textit{Insertion Sort} entre $k = 10^{5}$ e $k = 10^{6}$.}
  \label{fig:range}
\end{figure}

\textbf{Algoritmos de comparação.} O \textit{Insertion Sort}
(${\approx}\,2{,}0 \times 10^{-3}$\,s) e o \textit{Merge Sort}
(${\approx}\,3{,}0 \times 10^{-4}$\,s) apresentam linhas horizontais ao
longo de toda a faixa de~$k$, como previsto, pois sua complexidade depende
apenas de~$n$.

\textbf{Radix Sort.} Exibe uma leve elevação com $k$, pois $k$ maior implica
mais dígitos ($d = \lceil \log_{10} k \rceil$) e mais passadas: o tempo sobe
de ${\approx}\,6{,}2 \times 10^{-5}$\,s ($k = 10^{3}$) para
${\approx}\,1{,}1 \times 10^{-4}$\,s ($k = 10^{6}$), apenas
${\approx}\,1{,}8$ vez enquanto $k$ cresce $10^{3}$ vezes.

\textbf{Counting Sort.} É o mais sensível a~$k$. Para $k = n = 10^{3}$, é o
mais rápido de todos (${\approx}\,2{,}5 \times 10^{-5}$\,s), cerca de 2,5
vezes mais que o \textit{Radix Sort} e 12 vezes mais que o \textit{Merge
Sort}: é o seu regime ideal, com complexidade efetiva $\mathcal{O}(2n)$.
Com $n$ fixo, o tempo escala linearmente com~$k$, conforme
$\mathcal{O}(n + k)$, e ocorrem três inversões:

\begin{enumerate}
  \item Entre $k = 10^{3}$ e $10^{4}$ ($k/n$ de $1$ a $10$), é ultrapassado
        pelo \textit{Radix Sort}: em $k = 10^{3}$ ainda é mais rápido
        (${\approx}\,2{,}5 \times 10^{-5}$\,s contra
        ${\approx}\,6{,}2 \times 10^{-5}$\,s), mas em $k = 10^{4}$ já não
        é (${\approx}\,8{,}3 \times 10^{-5}$\,s contra
        ${\approx}\,6{,}7 \times 10^{-5}$\,s).
  \item Entre $k = 10^{4}$ e $10^{5}$ ($k/n$ de $10$ a $100$), ultrapassa o
        \textit{Merge Sort}: em $k = 10^{4}$ ainda é mais rápido
        (${\approx}\,8{,}3 \times 10^{-5}$\,s contra
        ${\approx}\,3{,}0 \times 10^{-4}$\,s), mas em $k = 10^{5}$ já não
        é (${\approx}\,6{,}5 \times 10^{-4}$\,s contra
        ${\approx}\,3{,}0 \times 10^{-4}$\,s).
  \item Entre $k = 10^{5}$ e $10^{6}$ ($k/n$ de $100$ a $1\,000$), ultrapassa
        também o \textit{Insertion Sort}: em $k = 10^{5}$ ainda está abaixo
        (${\approx}\,6{,}5 \times 10^{-4}$\,s contra
        ${\approx}\,2{,}0 \times 10^{-3}$\,s), mas em $k = 10^{6}$ torna-se
        o mais lento de todos (${\approx}\,6{,}4 \times 10^{-3}$\,s contra
        ${\approx}\,2{,}1 \times 10^{-3}$\,s).
\end{enumerate}

Dois problemas surgem quando $k \gg n$: (i)~\textbf{degradação de
desempenho}, pois o algoritmo percorre e inicializa um vetor de contadores
de tamanho~$k$ mesmo que a maioria das posições fique em zero; e
(ii)~\textbf{consumo de memória auxiliar} $\mathcal{O}(k)$, que pode agravar
o tempo por pior aproveitamento de cache (hipótese não verificada neste
trabalho).

Portanto, a escolha do \textit{Counting Sort} deve ser guiada pela relação
$k/n$: quando $k \approx n$, ele é o mais rápido; já com $k/n = 10$ o
\textit{Radix Sort} o supera, e a partir de $k/n$ entre $10$ e $100$ o
\textit{Merge Sort} também passa a superá-lo; e, para $k$ maiores, o
\textit{Radix Sort} oferece o melhor equilíbrio entre sensibilidade a~$k$ e
desempenho absoluto, sendo mais rápido que o \textit{Merge Sort} em todos os
casos avaliados.

% -------------------------------------------------------------------
\subsection{Análise Espacial: Mapa de Calor do Counting Sort}
\label{subsec:heatmap}
% -------------------------------------------------------------------

A Figura~\ref{fig:heatmap} mostra o tempo do \textit{Counting Sort} em
função de $n$ (eixo X) e $k$ (eixo Y), em escala de cores logarítmica: do
amarelo claro (baixo tempo) ao vermelho-escuro (degradação severa).

\begin{figure}[!htbp]
  \centering
  \includegraphics[width=0.75\columnwidth]{Images/heatmap_counting_sort.png}
  \caption{Mapa de calor do tempo de execução do \textit{Counting Sort}
           em função de $n$ (eixo X) e $k$ (eixo Y), com escala de cores
           logarítmica. Tons amarelos indicam alto desempenho; tons
           vermelho-escuros indicam degradação severa. As células em
           branco correspondem a combinações de $n$ e $k$ que não foram
           testadas.}
  \label{fig:heatmap}
\end{figure}

\textbf{Diagonal de eficiência.} A região $k \approx n$ concentra os menores
tempos (tons amarelos), onde a complexidade efetiva é $\mathcal{O}(2n)$; o
menor tempo registrado é ${\approx}\,2{,}5 \times 10^{-5}$\,s, em
$n = k = 1\,000$.

\textbf{Degradação vertical.} Aumentando~$k$ com $n$ fixo, as cores passam
do amarelo ao vermelho-escuro, assinatura do custo $\mathcal{O}(k)$ da
varredura dos contadores. Para $n = 10^{5}$ e $k = 10^{8}$, o tempo foi de
${\approx}\,0{,}73$\,s, mais de 300 vezes o do ponto ideal da diagonal
(${\approx}\,2{,}3 \times 10^{-3}$\,s).

O mapa serve, assim, como ferramenta de decisão: dado um par $(n, k)$,
indica se o \textit{Counting Sort} opera em regime eficiente (região clara)
ou se \textit{Radix Sort} ou \textit{Merge Sort} são preferíveis (região
escura).

% -------------------------------------------------------------------
\subsection{Síntese Comparativa}
\label{subsec:sintese}
% -------------------------------------------------------------------

A Tabela~\ref{tab:sintese} consolida as complexidades teóricas e as
principais observações empíricas.

\begin{table}[!htbp]
  \centering
  \caption{Síntese das complexidades teóricas e do comportamento
           experimental dos algoritmos avaliados. $b$ é a base numérica
           usada pelo Radix Sort (aqui, $b = 10$).
           ``Cond.'' indica viabilidade condicional à relação $k/n$.}
  \label{tab:sintese}
  \begin{tabular}{lcccc}
    \toprule
    \textbf{Algoritmo} & \textbf{Complexidade} & \textbf{Espaço} &
    \textbf{Dep.\ de $k$} &
    \textbf{Viável p/ $n = 10^{5}$} \\
    \midrule
    Insertion Sort & $\mathcal{O}(n^{2})$    & $\mathcal{O}(1)$      & Não      & Não   \\
    Merge Sort     & $\mathcal{O}(n \log n)$ & $\mathcal{O}(n)$      & Não      & Sim   \\
    Radix Sort     & $\mathcal{O}(d \cdot n)$& $\mathcal{O}(n + b)$  & Indireta & Sim   \\
    Counting Sort  & $\mathcal{O}(n + k)$    & $\mathcal{O}(k)$      & Direta   & Cond. \\
    \bottomrule
  \end{tabular}
\end{table}

Não existe algoritmo universalmente superior: a escolha depende do perfil da
instância. Com $k \approx n$, o \textit{Counting Sort} é o mais
rápido (${\approx}\,2{,}3 \times 10^{-3}$\,s para $n = 10^{5}$). Quando $k$
cresce além de $n$, o \textit{Radix Sort} é o melhor compromisso: foi o mais
rápido na média sobre os intervalos e o segundo quando $k = n$. O
\textit{Merge Sort} é a melhor opção entre os algoritmos de comparação, com
desempenho independente de~$k$. O \textit{Insertion Sort} mostrou-se
inviável em escala: para $n = 10^{5}$, foi cerca de 500 vezes mais lento que
o \textit{Merge Sort} e quase quatro ordens de grandeza mais lento que o
\textit{Counting Sort} com $k = n$.

% ============================================================
\section{Conclusão}
\label{sec:conclusao}
% ============================================================

Este trabalho implementou e comparou experimentalmente o Insertion Sort, o
Merge Sort, o Counting Sort e o Radix Sort sobre vetores de inteiros
aleatórios, variando o tamanho $n$ da entrada e o intervalo de valores $k$.
Os resultados confirmaram as análises teóricas: o Insertion Sort apresentou
crescimento quadrático e levou cerca de $19{,}7$\,s para ordenar
$10^{5}$ elementos, tornando-se inviável em entradas de maior escala; o
Merge Sort manteve comportamento $\Theta(n \log n)$ e tempo independente de
$k$; e os algoritmos lineares exibiram crescimento próximo de linear
quando $k = \mathcal{O}(n)$.

A principal conclusão é que o parâmetro $k$ determina o desempenho dos
algoritmos lineares. O Counting Sort foi o mais rápido quando $k = n$
($\approx 2{,}3 \times 10^{-3}$\,s para $n = 10^{5}$), mas degradou-se
rapidamente à medida que $k$ superou $n$: já com $k = 10n$ foi superado
pelo Radix Sort, na média sobre todos os intervalos testados passou a ser
mais lento que o Merge Sort, e para $n = 1\,000$ e $k = 10^{6}$ chegou a ser
mais lento que o próprio Insertion Sort. Além do tempo, seu uso de memória
auxiliar $\mathcal{O}(k)$ reforça a necessidade de cautela quando o
intervalo de valores é grande. O Radix Sort, por sua vez, mostrou-se pouco
sensível a $k$, com variação de apenas ${\approx}\,1{,}8$ vez quando $k$
cresceu $10^{3}$ vezes, e foi o mais rápido na média sobre os intervalos
testados ($\approx 1{,}1 \times 10^{-2}$\,s para $n = 10^{5}$), o que o
torna a alternativa mais robusta entre os algoritmos lineares quando o
intervalo dos dados não é conhecido.

Conclui-se, portanto, que não existe um algoritmo universalmente superior,
e que a escolha deve considerar o perfil da instância, em particular a
relação $k/n$: o Counting Sort é preferível quando $k \approx n$; o
Radix Sort, quando $k$ é maior ou desconhecido; o Merge Sort oferece uma
garantia $\mathcal{O}(n \log n)$ independente de $k$ entre os algoritmos de
comparação; e o Insertion Sort fica restrito a entradas pequenas ou quase
ordenadas.

Os resultados valem para as condições avaliadas: vetores de inteiros com
distribuição uniforme, uma única máquina e uma única configuração de
compilação (sem otimizações), de modo que os tempos absolutos podem diferir
em código compilado com \texttt{-O2} ou \texttt{-O3}. Como trabalhos futuros, sugere-se avaliar outras distribuições
de dados (como vetores quase ordenados ou com muitos valores repetidos),
outras bases no Radix Sort e a comparação com algoritmos como Quick Sort e
Heap Sort.

\bibliographystyle{sbc}
\bibliography{sbc-template}

\end{document}